\begin{abstract}
Although large-scale unsupervised language models (LMs) acquire extensive world knowledge and some reasoning abilities, exerting precise control over their outputs remains challenging owing to the inherently unsupervised nature of their training paradigms. Current approaches to enhance such control typically involve gathering human feedback in the form of pairwise preference data and subsequently adapting the LM via fine-tuning so that its outputs better reflect human judgements, most commonly through reinforcement learning from human feedback (RLHF). However, RLHF tends to be both intricate and prone to instability, necessitating an initial stage in which a reward model is trained to capture human preferences, followed by reinforcement learning to optimize the LM according to this learned reward, all the while seeking to limit divergence from the base model. This work presents a novel reward model parameterization for RLHF, which facilitates closed-form derivation of the associated optimal policy and thereby allows the conventional RLHF objective to be addressed using only a straightforward classification loss. The proposed method, termed \textit{Direct Preference Optimization} (DPO), achieves robust, effective, and computationally efficient alignment, circumventing the need for language model sampling during fine-tuning and minimizing the requirement for elaborate hyperparameter tuning. Empirically, DPO enables LMs to be fine-tuned in a manner that aligns with human preferences as effectively as, or more effectively than, prevalent techniques. Importantly, DPO surpasses PPO-based RLHF in steering generative sentiment, and matches or improves upon existing systems in both summarization and single-turn dialogue benchmarks, while also offering significant simplicity in implementation and training.
\end{abstract}

\section{Introduction}
Unsupervised language models (LMs) trained on vast corpora have demonstrated unexpectedly sophisticated abilities~\citep{chowdhery2022palm, brown2020language, touvron2023llama, bubeck2023sparks}. Nevertheless, the training data for these models is generated by humans whose objectives, priorities, and expertise vary widely. Some of these goals or skills may be inappropriate or undesirable for imitation; for instance, while it is beneficial for an AI coding assistant to \textit{recognize} frequent programming errors to provide corrections, ideally, the model’s code generation should be skewed toward reproducing the (often uncommon) high-caliber programming found in the data. In the same vein, it is useful for a language model to \textit{know about} a misconception held by half the population, yet undesirable for it to present that misconception as truth in half of the relevant user interactions. Thus, the ability to select and enforce the model’s \emph{intended behaviors and outputs} from its extensive \textit{repertoire of knowledge and skills} is essential for constructing AI systems that are safe, effective, and manageable~\citep{ouyang2022training}. Whereas prevailing approaches primarily align LMs with human values using reinforcement learning (RL), this work demonstrates that the RL-based objective employed by current strategies can in fact be directly solved via a straightforward binary cross-entropy loss, thereby streamlining the process of preference learning.

Generally, established methodologies incorporate desired traits into language models by leveraging carefully selected human preference datasets that encapsulate behaviors deemed safe and beneficial by people. This phase of preference learning succeeds a large-scale unsupervised pre-training process over extensive textual datasets. Although the simplest route for preference acquisition involves supervised fine-tuning on exemplary human-crafted responses, the most effective techniques to date use reinforcement learning from human (or artificial) feedback (RLHF/RLAIF; \citep{christiano2017deep,bai2022constitutional}). In RLHF, a reward model is trained on annotated human preferences and used to guide the LM policy through RL, aiming to produce outputs that maximize this reward without diverging significantly from the pretrained model. While RLHF has yielded models with remarkable performance in dialogue and programming tasks, the RLHF workflow remains much more involved than standard supervised learning, necessitating the training of multiple language models and policy sampling within the training loop—actions that substantially increase computational requirements.

In this work, we present a method to train language models directly based on human preference data, circumventing the need for explicit reward modeling or reinforcement learning frameworks. We introduce \textit{Direct Preference Optimization (DPO)}, an algorithm that, while implicitly maximizing the same reward objective constrained by KL-divergence as conventional RLHF techniques, features a streamlined implementation and straightforward training process. Conceptually, DPO modifies the model by increasing the log likelihood of preferred outputs relative to less favored ones, and uniquely incorporates a dynamic, sample-specific importance weight. This addition mitigates the model degradation observed when using a simplistic probability ratio approach. DPO, akin to prior algorithms, leverages a theoretical preference structure (for instance, the Bradley-Terry model; \cite{bradley1952rankanalysis}) to evaluate the agreement between reward functions and collected preference labels. Distinctly, rather than employing the preference model to establish a loss for reward model training before subsequent policy optimization, DPO reparameterizes the preference loss to directly depend on the policy itself. Utilizing collections of human-annotated preferences regarding model outputs, DPO optimizes the policy with a basic binary cross entropy loss, thereby attaining the optimal policy corresponding to an implicit reward landscape tailored to the preference dataset.

The central contribution of this study is the Direct Preference Optimization (DPO) algorithm, which dispenses with reinforcement learning in favor of a more straightforward preference-driven approach to language model training. Empirical evaluations reveal that DPO matches or surpasses the performance of established approaches, such as PPO-based RLHF, across various preference-learning tasks—including sentiment adjustment, summarization, and conversational modeling—when applied to language models containing up to 6 billion parameters.

\section{Related Work}

As self-supervised language models have grown in scale, their ability to perform various tasks with zero-shot \citep{radford2019language} or few-shot prompting \citep{gpt3,megatron,chowdhery2022palm} has improved. Nevertheless, augmenting these models with fine-tuning on curated datasets consisting of instructions and human-generated completions has proven to substantially enhance their effectiveness on downstream benchmarks and their alignment with user expectations \citep{mishra-etal-2022-cross,sanh2022multitask,chung2022scaling,thoppilan2022lamda}. This approach, commonly referred to as `instruction-tuning,’ not only expands LLMs' capacity to handle instructions beyond the scope of their fine-tuning set but also generally improves their practical utility \citep{chung2022scaling}. While instruction tuning has yielded notable gains, collecting \textit{relative} human assessments of response quality is often more scalable than soliciting expert-annotated demonstrations. Consequently, more recent research has explored the fine-tuning of LLMs using human preference data, yielding improvements across a range of tasks such as translation \citep{kreutzer-etal-2018-reliability}, summarization \citep{stiennon2022learning,ziegler2020finetuning}, story generation \citep{ziegler2020finetuning}, and following instructions \citep{ouyang2022training,ramamurthy2023is}. These strategies typically involve first training a reward model, based on neural networks, that predicts human preferences using frameworks like the Bradley-Terry model \citep{bradley1952rankanalysis}, followed by fine-tuning the language model to maximize this reward via reinforcement learning techniques—most notably REINFORCE \citep{williams1992reinforce}, proximal policy optimization (PPO; \cite{schulman2017proximal}), or related methods \citep{ramamurthy2023is}. Related research streams employ LLMs already tuned to follow instructions and accept human feedback to autonomously create additional synthetic preference datasets targeting criteria such as safety or harmlessness \citep{bai2022constitutional}, guided only by weak supervision in the form of textual guidelines for annotations. This intersection of methods bridges earlier work on reinforcement learning-based training of language models across diverse objectives~\citep{Ranzato2015SequenceLT,paulus2018a,wu2018learning} with general approaches for learning from human preference data \citep{christiano2017deep,kupcsik2018learning}. Although leveraging relative human preferences presents an attractive alternative, the process of applying reinforcement learning to adapt large language models remains a substantial practical hurdle; the present work introduces a principled method to directly optimize for relative preferences without necessitating reinforcement learning.

Beyond language-related applications, policy learning from preferences has been extensively examined within both bandit and reinforcement learning frameworks, giving rise to various methodologies. Contextual bandit approaches that utilize rankings or preferences in lieu of explicit rewards are commonly referred to as contextual dueling bandits (CDB; \cite{yue2012karmed,dudik2015contextual}). When absolute reward signals are unavailable, the theoretical treatment of CDBs replaces the optimal policy with the concept of a \textit{von Neumann winner}, defined as a policy whose expected success rate against any alternative policy meets or exceeds 50\% \citep{dudik2015contextual}. A notable distinction is that, in the CDB formulation, preference annotations are typically acquired in an online manner. By contrast, in the context of learning from human preferences, the process often relies on a fixed, offline set of action pairs annotated with preference information \citep{yan2022human}. Likewise, \textit{preference-based RL} (PbRL) aims to learn using binary preference signals, which are outputs of an unknown ‘scoring’ function instead of direct rewards \citep{BusaFekete2014,ruiz2023dueling}. PbRL comprises several algorithmic strategies—including those capable of reusing off-policy preference datasets—but these generally involve an explicit phase for estimating the underlying reward (scoring) function before proceeding to optimize it \citep{jain2013learning,BusaFekete2014,christiano2017deep,sadigh2017active,kupcsik2018learning}. Contrasting with this two-step procedure, our contribution is a unified approach that learns policies directly optimized with respect to observed preferences.

\section{Preliminaries}\label{section:prelims}

Here, we summarize the RLHF methodology presented in \citeauthor{ziegler2020finetuning} and further extended in works such as \citep{stiennon2022learning, bai2022training, ouyang2022training}. The standard RLHF pipeline is divided into three distinct phases: (1) supervised fine-tuning (SFT); (2) sampling for preference annotation and reward modeling; and (3) policy optimization via reinforcement learning.

\textbf{SFT}: The RLHF process begins by taking a pre-trained language model and applying supervised fine-tuning using curated, task-relevant datasets (e.g., for dialogue, summarization), resulting in a policy denoted as $\pi^\text{SFT}$. 

\textbf{Reward Modeling Phase}: During this stage, the SFT policy generates responses to prompts $x$ in the form of answer pairs $(y_1, y_2)\sim \pi^\text{SFT}(y \mid x)$. These pairs are then evaluated by human annotators, who express a preference by selecting one completion over the other—indicated by $y_w\succ y_l \mid x$, where $y_w$ and $y_l$ signify the more and less preferred answers, respectively. These expressed preferences are treated as outputs from an unobserved, latent reward function $r^*(y, x)$. While a variety of statistical models can be employed to capture these preferences, the Bradley-Terry (BT) model \cite{bradley1952rankanalysis} is frequently utilized; more sophisticated ranking models such as Plackett-Luce \citep{plackett1975analysis, luce2012individual} are also applicable if multiple ranked responses are collected. The BT formulation posits the human preference probability $p^*$ in the following way:

\begin{equation}\label{eq:bradley-terry}
    p^*(y_1\succ y_2 \mid x)=\frac{\exp\left(r^*(x, y_1)\right)}{\exp\left(r^*(x, y_1)\right) + \exp\left(r^*(x, y_2)\right)}.
\end{equation}

Given a fixed collection of pairwise preference data $\mathcal{D}=\bigl\{x^{(i)}, y_w^{(i)}, y_l^{(i)}\bigr\}_{i=1}^N$ drawn from $p^*$, one can introduce a parameterized reward function $r_{\phi}(x, y)$ and optimize its parameters via maximum likelihood estimation. Framing the task as binary classification leads to the following negative log-likelihood objective:

\begin{equation}\label{eq:reward_model}
    \mathcal{L}_R(r_{\phi}, \mathcal{D}) = -\mathbb{E}_{(x, y_w, y_l)\sim \mathcal{D}}\bigl[\log \sigma(r_{\phi}(x, y_w)- r_{\phi}(x, y_l))\bigr]
\end{equation}

where $\sigma$ denotes the logistic sigmoid function. For large language models, $r_{\phi}(x, y)$ is commonly initialized with the SFT model $\pi^\text{SFT}(y \mid x)$, appending a linear head atop the last transformer block to yield a scalar reward score~\cite{ziegler2020finetuning}. To reduce reward variance, it is standard practice to normalize the outputs such that $\mathbb{E}_{x,y\sim \mathcal{D}}\left[r_\phi(x, y)\right] = 0$ across $x$.

\textbf{RL Fine-Tuning Phase}: In the subsequent reinforcement learning stage, this learned reward is used as supervisory feedback for model training. As in earlier work~\citep{jaques2017sequence, jaques2020human}, policy optimization seeks to solve

\begin{equation}\label{eq:RL}
\max_{\pi_{\theta}}  \mathbb{E}_{x\sim \mathcal{D}, y\sim \pi_{\theta}(y \mid x)}\bigl[r_{\phi}(x, y)\bigr] - \beta\mathbb{D}_{\textrm{KL}}\bigl[\pi_{\theta}(y\mid x)\mid \mid \pi_\text{ref}(y\mid x)\bigr],
\end{equation}

where the coefficient $\beta$ regulates divergence from a reference policy $\pi_\text{ref}$, typically set as the initial SFT model $\pi^\text{SFT}$. 
Typically, the policy $\pi_\theta$ is also initialized to $\pi^\text{SFT}$. This regularization is critical: it constrains the updated model to remain near the regime where reward evaluations are meaningful, supports diversity in generations, and mitigates collapse onto a few reward-maximizing outputs. Given the non-differentiable, discrete nature of text outputs, this optimization is not amenable to gradient descent and is instead addressed with reinforcement learning techniques. Conventional strategies~\citep{ziegler2020finetuning, stiennon2022learning, bai2022training, ouyang2022training} formulate the reward as ${r(x, y) = r_{\phi}(x, y) -\beta (\log \pi_{\theta}(y\mid x) - \log \pi_\text{ref}(y\mid x))}$, maximizing it with algorithms such as PPO~\cite{schulman2017proximal}.

\section{Direct Preference Optimization}\label{sec:DPO}

In response to the complexity inherent in using reinforcement learning for large-scale language model fine-tuning, our objective is to introduce a more straightforward preference-based policy optimization framework. Departing from earlier RLHF techniques that decouple reward modeling and RL-based optimization, our strategy hinges on choosing a reward parameterization that permits a closed-form solution for the optimal policy, entirely bypassing the need for an explicit RL loop. We now provide a detailed explanation: our central idea is to exploit an analytical correspondence between reward functions and their induced optimal policies, effectively reformulating a reward-based loss into a policy-based objective through a variable transformation. This method precludes the necessity for a separately trained reward model, yet continues to align with probabilistic human preference models, such as the Bradley-Terry formulation. As a result, the policy itself encodes both the language generation mechanism and the implicit reward structure.

\textbf{DPO Objective Derivation.} We begin from the standard RL framework used in earlier studies, Eq.~\ref{eq:RL}, considering an arbitrary reward function $r$. Consistent with previous research~\citep{peters2007reinforcement, peng2019advantage, korbak2022reinforcement, go2023aligning}, it is evident that the solution optimizing the KL-regularized reward maximization problem in Eq.~\ref{eq:RL} is characterized by:

\begin{equation}\label{eq:op_policy}
    \pi_r(y\mid x) = \frac{1}{Z(x)}\pi_\text{ref}(y\mid x)\exp\left(\frac{1}{\beta}r(x, y)\right),
\end{equation}

Here, the partition function $Z(x) =\sum_{y}\pi_\text{ref}(y\mid x)\exp\left(\frac{1}{\beta}r(x, y)\right)$ ensures proper normalization. Further mathematical details are available in Appendix~\ref{app:derivation1}. In practice, even using an MLE-optimized approximation $r_\phi$ of the true reward function $r^*$ does not make evaluating $Z(x)$ computationally efficient \citep{korbak2022reinforcement, go2023aligning}, thereby limiting the practical application of this form. Nevertheless, Eq.~\ref{eq:op_policy} can be reformulated so the reward function is described in terms of the corresponding optimal policy $\pi_r$, the baseline policy $\pi_\text{ref}$, and the unknown normalization constant $Z(\cdot)$. Taking the logarithm on both sides of Eq.~\ref{eq:op_policy} and rearranging yields:

\begin{equation}\label{eq:main_eq}
    r(x,y) =\beta \log \frac{\pi_r(y\mid x)}{\pi_\text{ref}(y\mid x)} + \beta \log Z(x).
\end{equation}

This formulation can be extended to the ground-truth reward $r^*$ and its optimal policy $\pi^*$. Crucially, the Bradley-Terry model relies solely on differences between the rewards for two options: specifically, ${p^*(y_1 \succ y_2 \mid x) = \sigma(r^*(x, y_1) - r^*(x, y_2))}$. Substituting the representation from Eq.~\ref{eq:main_eq} for $r^*(x,y)$ into the preference formulation in Eq.~\ref{eq:bradley-terry} removes the partition function term, permitting an expression for the probability of human preference that involves only the optimal policy $\pi^*$ and the reference policy $\pi_\text{ref}$. Consequently, the policy $\pi^*$ optimal for RLHF under the Bradley-Terry framework fulfills the following preference probability relation:

\begin{equation}\label{eq:objective}
    p^*(y_1\succ y_2 \mid x)=\frac{1}{1 + \exp\left(\beta \log \frac{\pi^*(y_2\mid x)}{\pi_\text{ref}(y_2\mid x)} - \beta \log \frac{\pi^*(y_1\mid x)}{\pi_\text{ref}(y_1\mid x)}\right)}
\end{equation}

A detailed step-by-step derivation can be found in Appendix~\ref{app:derivation2}. Although Eq.~\ref{eq:objective} leverages the Bradley-Terry structure, similar derivations extend to the broader class of Plackett-Luce models~\citep{plackett1975analysis, luce2012individual} as elaborated in Appendix~\ref{app:plackett_luce_models}.

With an explicit form for the probability of observed human preferences based solely on the optimal policy (and independent of the reward function parameterization), it is possible to establish a likelihood-based objective to train a parameterized policy $\pi_\theta$. Drawing a parallel with the approach for reward modeling (as in Eq.~\ref{eq:reward_model}), we arrive at the following policy optimization objective:

\begin{equation}\label{eq:optimum_model}
    \mathcal{L}_\text{DPO}(\pi_{\theta}; \pi_\text{ref}) = -\mathbb{E}_{(x, y_w, y_l)\sim \mathcal{D}}\left[\log \sigma \left(\beta \log \frac{\pi_{\theta}(y_w\mid x)}{\pi_\text{ref

In this formulation, we model the implicit reward through an alternative parameterization, wherein the optimal policy corresponds directly to $\pi_\theta$. Additionally, because this approach is analogous to learning a reparametrized Bradley-Terry model, it inherits desirable theoretical guarantees, including consistency provided that the preference data satisfy appropriate assumptions \cite{bong2022generalized}. We elaborate on further theoretical aspects of DPO and comparisons to related methodologies in Section~\ref{sec:theory}.

\textbf{How does the DPO update function?} To gain deeper mechanistic insight into DPO, it is instructive to inspect the gradient of the loss $\mathcal{L}_\text{DPO}$. The derivative with respect to the model parameters $\theta$ is given by:

\begin{multline*}\label{eq:gradient}
    \nabla_\theta \mathcal{L}_\text{DPO}(\pi_\theta;\pi_\text{ref}) = \\ -\beta\mathbb{E}_{(x, y_w, y_l) \sim \mathcal{D}} \bigg[\underbrace{\sigma(\hat{r}_\theta(x, y_l) - \hat{r}_\theta (x, y_w))}_\text{places more emphasis when the reward prediction is incorrect}\bigg[\underbrace{\nabla_\theta\log \pi(y_w \mid x)}_\text{encourage $y_w$} - \underbrace{\nabla_\theta\log\pi(y_l \mid x)}_\text{discourage $y_l$}\bigg]\bigg],
\end{multline*}

Here, the reward $\hat{r}_\theta(x, y)$ is defined as $\beta \log \frac{\pi_\theta(y \mid x)}{\pi_\text{ref}(y \mid x)}$, implicitly determined by the current policy $\pi_\theta$ and a reference policy $\pi_\text{ref}$ (see additional discussion in Section~\ref{sec:theory}). Conceptually, the gradient step encourages higher probability for the chosen or preferred response $y_w$ and suppresses the likelihood of the less-preferred response $y_l$. Crucially, each sample is weighted according to the extent to which the implicit reward model $\hat{r}_\theta$ erroneously rates $y_l$ above $y_w$, scaled by $\beta$, thus reflecting the severity of misordering by the reward model as well as the strength of the KL regularization. Empirical findings highlight the importance of this weighting mechanism, as omitting it—by not applying the weighting factor—leads to pathological model behaviors, as evidenced in Appendix Table~\ref{tab:unlikelihood_generations}.

\textbf{Overview of DPO.} The typical DPO workflow proceeds in two main steps: First, for each input prompt $x$, generate candidate completions $y_1, y_2$ by sampling from $\pi_\text{ref}(\cdot \mid x)$, and annotate these completions with human preference judgments to assemble an offline preference dataset $\mathcal{D} = \{x^{(i)}, y_w^{(i)}, y_l^{(i)}\}_{i=1}^N$. Second, train the language model $\pi_\theta$ by minimizing $\mathcal{L}_\text{DPO}$, using the specified $\pi_\text{ref}$, preference data $\mathcal{D}$, and a chosen $\beta$ value.

In practical scenarios, it is desirable to make use of existing publicly available preference datasets instead of independently generating responses and collecting new preference annotations. As these datasets are typically produced by sampling from a supervised fine-tuned model $\pi^\text{SFT}$, we opt to set $\pi_\text{ref} = \pi^\text{SFT}$ whenever such a model is present. In cases where $\pi^\text{SFT}$ cannot be accessed, $\pi_\text{ref}$ is instead initialized by maximizing the likelihood over preferred completions, formally, ${\pi_\text{ref} = \argmax_{\pi}\mathbb{E}_{x, y_w \sim \mathcal{D}}\left[\log \pi(y_w \mid x)\right]}$. This initialization aims to alleviate the potential distribution gap between the unknown true reference distribution and the one represented by $\pi_\text{ref}$ in DPO. Additional implementation details and information on hyperparameter selection are provided in Appendix~\ref{app:implementation}.

\section{Theoretical Analysis of DPO} In this section, we deepen the understanding of DPO by offering theoretical insight, formal justification, and by elucidating how the strengths of DPO address shortcomings commonly encountered with actor-critic methods used in RLHF, including algorithms like PPO~\cite{schulman2017proximal}.

\label{sec:theory}

\subsection{Your Language Model Is Secretly a Reward Model} DPO manages to both avoid constructing an explicit reward function and sidestep reinforcement learning-based policy optimization, instead using a unified maximum likelihood framework. Notably, the loss objective given in Eq. \ref{eq:main_eq} aligns with the Bradley-Terry model when employing a reward structure of $r^*(x, y) = \beta \log\frac{\pi^*_\theta(y \mid x)}{\pi_\text{ref}(y \mid x)}$, such that optimizing the parametric policy model $\pi_{\theta}$ corresponds—up to a variable transformation—to fitting a reward model as in Eq. \ref{eq:reward_model}. In what follows, we develop the theoretical basis for this reparameterization, demonstrate its flexibility for representing all possible learned reward models, and prove that it permits exact recovery of the optimal policy. We begin our exposition by introducing a relevant equivalence relation between reward functions.

\begin{definition}
We define two reward functions $r(x, y)$ and $r'(x, y)$ to be equivalent if and only if there exists some function $f$ such that $r(x, y)-r'(x, y) = f(x)$.
\end{definition}

This definition immediately yields an equivalence relation, thereby dividing the collection of reward functions into equivalence classes. The following lemmas can be established:

\begin{lemma}\label{lemma:same_prefrence} Within the Plackett-Luce framework, including the Bradley-Terry model as a special case, any pair of reward functions belonging to the same equivalence class give rise to identical preference distributions.
\end{lemma}

\begin{lemma}\label{lemma:same_policy}
    For the constrained RL setting, any two reward functions from a single equivalence class result in the same optimal policy.
\end{lemma}

The proofs are direct and are postponed to Appendix \ref{app:lemma1}. The first lemma highlights a well-documented issue of under-identification inherent in the Plackett-Luce family of models \cite{plackett1975analysis}. This ambiguity necessitates additional identifiability constraints when seeking reliable MLE estimates for Eq.~\ref{eq:reward_model} \cite{bong2022generalized}. The second lemma asserts that all reward functions within the same equivalence class lead to an identical optimal policy. Thus, for our purposes, it suffices to identify any reward function within the corresponding optimal class. We formalize this in the following theorem, proven in Appendix~\ref{app:thm1}:

\begin{theorem}\label{thm:main}
    Assuming mild conditions, every equivalence class of reward functions compatible with the Plackett-Luce model (and specifically the Bradley-Terry model) admits a representation of the form ${r(x, y) = \beta \log \frac{\pi(y\mid x)}{\pi_\text{ref}(y\mid x)}}$ for some choice of model $\pi(y\mid x)$, with a fixed reference model $\pi_\text{ref}(y \mid x)$.
\end{theorem}

\begin{sproof}
    Take an arbitrary reward function $r(x, y)$, and denote its associated optimal policy by $\pi_r(y \mid x)$ as specified in Eq.~\ref{eq:op_policy}. We aim to demonstrate that any member of the equivalence class containing $r$ can be captured through the proposed reparameterization. Define the projection operator $f$ as follows:
\begin{equation}
    f(r; \pi_\text{ref}, \beta)(x, y) = r(x, y) - \beta\log\sum_{y}\pi_\text{ref}(y\mid x)\exp\left(\frac{1}{\beta}r(x, y)\right)
\end{equation}
Here, $f$ adjusts the reward by subtracting the log-partition function of $\pi_r$ (a function solely of $x$), ensuring that $f(r; \pi_\text{ref}, \beta)(x, y)$ remains in the same equivalence class as $r(x, y)$. Substituting the expression for $r$ from Eq.~\ref{eq:main_eq} (which is applicable to any reward function), it follows that $f(r; \pi_\text{ref}, \beta)(x, y) = \beta \log \frac{\pi_r(y\mid x)}{\pi_\text{ref}(y\mid x)}$. Thus, the projection $f$ yields a reward function of the desired form within the equivalence class of $r$, establishing that our chosen reparameterization encompasses all generality needed for the reward model.
\end{sproof}

An alternative perspective on Theorem~\ref{thm:main} is that it uniquely determines the reward function, within each equivalence class, chosen by the DPO reparameterization. Specifically, it selects the reward function that fulfills:

\begin{equation}\label{eq:lag_p}
     \sum_{y}\underbrace{\pi_\text{ref}(y\mid x)\exp\left(\frac{1}{\beta}r(x, y)\right)}_{=\pi(y\mid x)\text{, by Thm.~\ref{thm:main} reparam.}} = 1,
\end{equation}

meaning that $\pi(y\mid x)$ constitutes a valid probability distribution (the probabilities are non-negative and sum to 1). Observing Eq.~\ref{eq:op_policy}, we recognize that Eq.~\ref{eq:lag_p} describes the partition function corresponding to the optimal policy dictated by the reward function $r(x, y)$. The central innovation in the DPO approach is to enforce constraints on the inherently under-specified Plackett-Luce (and, by extension, Bradley-Terry) class of preference models, maintaining the expressiveness of the reward model class, while also ensuring that the optimal policy in Eq.~\ref{eq:op_policy} remains analytically solvable for all prompt instances $x$.

\subsection{Instability of Actor-Critic Algorithms} 

Our formalism can further shed light on the causes of instability in common actor-critic algorithms, such as PPO, which are standard in RLHF pipelines. We focus on the RL fine-tuning component as introduced in Section \ref{section:prelims}. This analysis draws on the control as inference paradigm \cite{levine2018reinforcement}, in relation to the constrained RL scenario described by \ref{eq:RL}. Let us suppose the use of a parametric model $\pi_{\theta}(y\mid x)$, with the training goal to minimize $\mathbb{D}_{\text{KL}}[\pi_{\theta}(y|x) \mid\mid \pi^*(y\mid x)]$, where $\pi^*$ is the optimal policy from Eq.~\ref{eq:optimum_model} based on the reward $r_{\phi}(y, x)$. Algebraic manipulation yields the following objective for optimization:

\begin{equation}\label{eq:AC}
    \max_{\pi_{\theta}}\mathbb{E}_{\pi_{\theta}(y\mid x)}\bigg[\underbrace{r_{\phi}(x, y) -\beta\log\sum_{y}\pi_\text{ref}(y\mid x)\exp\left(\frac{1}{\beta}r_{\phi}(x, y)\right)}_{f(r_{\phi}, \pi_\text{ref}, \beta)} - \underbrace{\beta\log\frac{\pi_{\theta}(y\mid x)}{\pi_\text{ref}(y\mid x)}}_{\text{KL}}\bigg]
\end{equation}

This objective has also been the focus of optimization in earlier studies 
\citep{ziegler2020finetuning, stiennon2022learning, bai2022training, ouyang2022training}, where the reward function $r_{\phi}$ utilized the DPO-equivalent reward structure. Within this context, the normalization component in $f(r_{\phi}, \pi_\text{ref}, \beta)$ can be viewed as the soft value function corresponding to the reference policy $\pi_\text{ref}$. Although the normalization term does not influence the optimal policy, omitting it may lead to a high-variance policy gradient, thereby hindering the stability of the learning process. One approach to address normalization is to introduce a learned value function, though this presents optimization challenges of its own. Previous research has often adopted a reward normalization method based on a human completion baseline, essentially using a Monte Carlo estimate with a single sample to approximate the normalization factor. The DPO formulation, by contrast, provides a reward function inherently independent of such baseline requirements.

\section{Experiments}
This section presents an empirical investigation into how effectively DPO can learn policies from preference signals. Initially, we conduct experiments within a controlled environment for text generation to explore DPO's efficiency in negotiating the trade-off between reward maximization and KL-divergence minimization relative to the reference policy, benchmarking its performance against established algorithms such as PPO for preference optimization. Subsequently, we apply DPO to larger-scale models and tackle more complex RLHF scenarios, including tasks like summarization and dialogue. Our observations indicate that DPO often matches or surpasses robust baselines—such as RLHF leveraging PPO or the strategy of selecting the best trajectory out of $N$ samples based on a learned reward—without substantial hyperparameter tuning. The details of our experimental methodology are outlined before the presentation of results; further specifics are available in Appendix~\ref{app:exp_details}.

\textbf{Tasks.} Our empirical study encompasses three distinct open-ended text generation tasks. In all cases, the algorithms optimize a policy based on a preference dataset $\mathcal{D}=\bigl\{x^{(i)}, y_w^{(i)}, y_l^{(i)}\bigr\}_{i=1}^N$. In the \textbf{controlled sentiment generation} task, $x$ represents an initial segment of a movie review sourced from the IMDb dataset \cite{maas-EtAl:2011:ACL-HLT2011}, with the requirement that the policy generates a continuation $y$ exhibiting positive sentiment. To systematically evaluate performance, this experiment constructs pairs of generated continuations and applies a pre-trained sentiment classifier to assign preferences, selecting those for which $p(\text{positive}\mid x,y_w)>p(\text{positive}\mid x,y_l)$. The SFT baseline is obtained by further training GPT-2-large on positive reviews from the IMDb training split until performance plateaus (additional details are provided in App~\ref{app:sentiment_details}). In the \textbf{summarization} setting, $x$ corresponds to a Reddit forum post and the policy is tasked with producing a summary $y$ highlighting the key points. In accordance with previous approaches, we utilize the Reddit TL;DR summarization dataset \citep{volske-etal-2017-tl} augmented with human preferences collected by \citeauthor{stiennon2022learning}. We initialize with an SFT model refined on human-authored post summaries\footnote{\url{https://huggingface.co/CarperAI/openai_summarize_tldr_sft}} and leverage the TRLX framework \citep{leandro_von_werra_2023_7790115} to facilitate RLHF training. The set of human preferences stems from annotations on outputs of a comparable, but not identical, SFT model. For the \textbf{single-turn dialogue} task, $x$ takes the form of a user prompt, ranging from scientific inquiries to personal requests, with the goal being to generate an informative and engaging response $y$. This uses the Anthropic Helpful and Harmless dialogue dataset \citep{bai2022training}, which contains 170,000 conversational exchanges between humans and an AI assistant; each concludes with two assistant responses generated by a large language model and a label indicating which the human raters preferred. Because an SFT baseline is not pre-trained in this context, we train an off-the-shelf language model on preferred responses alone to derive the SFT model for this task.

\textbf{Evaluation.} We employ two distinct strategies for evaluation in our experiments. To assess how well each algorithm meets the constrained reward maximization objective in a controlled sentiment generation context, we measure each algorithm’s performance along the Pareto frontier of achieved reward and KL-divergence from the reference policy. This is feasible in this setting because the true reward function (derived from a sentiment classifier) is accessible. Conversely, in practical scenarios where the ground-truth reward function is unavailable, we assess algorithm performance based on the \textit{win rate} when compared to a baseline policy. Here, GPT-4 acts as a surrogate for human evaluators, judging summary quality and response helpfulness in summarization and single-turn dialogue tasks, respectively. For summarization, the baseline is the reference summaries from the test set; for dialogue, it is the preferred response from the test data. Prior research has indicated that language models can serve as more reliable automated evaluators than traditional metrics \citep{Chen2023ExploringTU}. Nonetheless, to validate our reliance on GPT-4 for evaluation, we conduct a human study as described in Sec.~\ref{sec:human-judgments}. Our findings show that GPT-4's evaluations exhibit strong agreement with those made by humans, often matching or exceeding the inter-annotator agreement levels among human raters.

\textbf{Methods.} Alongside DPO, we benchmark several established methods for aligning language model outputs with human preferences. For the summarization task, we utilize zero-shot prompting with \textbf{GPT-J} \citep{gpt-j}, while in the dialogue task, we apply 2-shot prompting with \textbf{Pythia-2.8B} \citep{biderman2023pythia}. Additionally, we include the \textbf{SFT} model and the \textbf{Preferred-FT} variant, which is trained via supervised fine-tuning on selected completions $y_w$, sourced from either the SFT model (for controlled sentiment and summarization) or a generic LM (for single-turn dialogue). Another technique, classified as pseudo-supervised, is the \textbf{Unlikelihood} method~\citep{welleck2019neural}, which aims to maximize the likelihood of $y_w$ and explicitly \textit{minimize} the likelihood of $y_l$, modulated by an optional unlikelihood coefficient $\alpha \in [0,1]$. We also evaluate \textbf{PPO} \citep{schulman2017proximal}, where the policy is optimized using a reward model inferred from preference annotations, and \textbf{PPO-GT}, an oracle approach utilizing the authentic reward function in the controlled sentiment scenario. For the sentiment generation experiments, two versions of PPO-GT are compared: an out-of-the-box implementation \cite{leandro_von_werra_2023_7790115} and a tailored version incorporating reward normalization and hyperparameter tuning to boost effectiveness (the same adjustments are adopted for the standard PPO when optimizing on learned rewards). Finally, we test the \textbf{Best of $N$} baseline, which involves sampling $N$ outputs from the SFT model (or Preferred-FT in dialogue), selecting the highest-reward response according to a reward model trained on the preference dataset. While this method offers strong performance by decoupling reward model quality from PPO-based optimization, it proves computationally expensive, since it necessitates sampling $N$ responses for every query during evaluation.

\subsection{How well can DPO optimize the RLHF objective?}

In conventional RLHF approaches, the reward maximization objective is constrained by KL divergence to prevent the learned policy from straying excessively from the reference policy, thereby enforcing a trade-off between maximizing reward and maintaining policy similarity. As such, when comparing different methods, it is essential to evaluate both the achieved reward and the associated KL divergence; modest gains in reward that come at the cost of a significant KL increase may not be advantageous. Figure~\ref{fig:frontier-tldr-main} illustrates the reward-KL trade-off for several algorithms in a sentiment analysis context. For each method, we perform several training runs, each time selecting a distinct value for the conservativeness hyperparameter (target KL $\in\{3,6,9,12\}$ for PPO, $\beta \in \{0.05,0.1,1,5\}$, $\alpha\in\{0.05,0.1,0.5,1\}$ for unlikelihood, and varying random seeds for preferred-FT), leading to a total of 22 experimental configurations. Throughout training, after every 100 steps until reaching convergence, we assess each policy on a fixed collection of test prompts, reporting both the mean reward under the ground-truth reward function and the average sequence-level KL\footnote{The sum over per-timestep KL-divergences.} with respect to the reference policy $\text{KL}\left(\pi\mid \mid \pi_\text{ref}\right)$. The experiments demonstrate that DPO establishes a substantially more favorable frontier, consistently yielding higher rewards with relatively low KL divergence. This is significant for several reasons. Firstly, although DPO and PPO share the same underlying objective, DPO is markedly more sample-efficient, with its reward/KL trade-off universally surpassing that of PPO. Furthermore, DPO maintains this superior frontier even in scenarios where PPO is granted access to ground-truth rewards (PPO-GT).

\subsection{Can DPO scale to real preference datasets?}
\label{sec:dpo-real-datasets}
We proceed to assess how well DPO can be fine-tuned on practical tasks involving real-world preference data, specifically focusing on summarization and single-turn dialogue. In the case of summarization, prior studies have shown that commonly used automated metrics like ROUGE often correlate poorly with human judgments~\citep{stiennon2022learning}. Earlier research demonstrated that language models fine-tuned with PPO on datasets reflecting human preferences are capable of producing higher-quality summaries. To compare various methods, we generate samples on the test portion of the TL;DR summarization dataset and calculate the average win rate relative to the reference outputs from the test set. Sampling is conducted across a temperature range from 0.0 to 1.0 for all approaches, with the resulting win rates illustrated in Figure~\ref{fig:frontier-tldr-main} (right). The models DPO, PPO, and Preferred-FT are all fine-tuned starting from the same GPT-J SFT model\footnote{\url{https://huggingface.co/CarperAI/openai_summarize_tldr_sft}}. Our analysis reveals that DPO attains a win rate close to 61\% at a temperature setting of 0.0, surpassing PPO, which achieves roughly 57\% at its own optimal temperature of 0.0. Furthermore, DPO outperforms the strongest best-of-$N$ baseline in terms of maximum win rate. It is important to note that we did not carry out substantial hyperparameter optimization for DPO's $\beta$, suggesting that the current results may not reflect the full capabilities of DPO. We also observe that DPO is substantially less sensitive to sampling temperature than PPO, whose effectiveness can drop to that of the underlying GPT-J model as the temperature increases. In contrast, Preferred-FT offers negligible improvements over the SFT baseline. A direct comparison between DPO and PPO using human assessments is presented in Section~\ref{sec:human-judgments}, indicating that samples from DPO at a temperature of 0.25 were favored 58\% of the time when compared to those from PPO at the same setting.

For single-turn dialogue experiments, we assess various approaches on the subset of the Anthropic HH dataset \citep{bai2022training} test split consisting of one round of human-assistant exchange. In these experiments, GPT-4 is used as an evaluator, with the preferred test completions serving as the reference to determine win rates among the different techniques. Since there is no established SFT model for this scenario, our pipeline begins with the pre-trained Pythia-2.8B, upon which we apply Preferred-FT to fine-tune a reference model on the selected completions, ensuring model outputs are in-distribution. DPO training is then performed on top of this setup. For comparison, we also consider the Best of 128 Preferred-FT completions (noting that performance for the Best of $N$ baseline saturates at $N=128$ on this task; see Appendix Figure~\ref{fig:best-of-n}), as well as a 2-shot prompt configuration using the base Pythia-2.8B, observing that DPO achieves comparable or superior win rates across optimal temperature selections for each method. In addition, we include evaluation of an RLHF agent fine-tuned using PPO on the Anthropic HH dataset\footnote{\url{https://huggingface.co/reciprocate/ppo_hh_pythia-6B}} from a notable implementation\footnote{\url{https://github.com/CarperAI/trlx/tree/main/examples/hh}}, but are unable to identify any prompt or sampling temperature yielding better outcomes than the Pythia-2.8B baseline. Drawing from findings in TL;DR and acknowledging that both PPO and Best of $N$ optimize towards the same reward objective, we interpret the Best of 128 baseline as an approximate stand-in for PPO-level capability. In summary, DPO stands out as the sole method that is both computationally efficient and effective at surpassing preferred completion baselines in the Anthropic HH dataset, offering performance comparable to or exceeding the resource-intensive Best of 128 baseline. Moreover, Figure~\ref{fig:dialogue-main} illustrates that DPO achieves near-optimal performance with rapid convergence.

\subsection{Generalization to a new input distribution}

To gain further insights into how PPO and DPO policies behave under shifts in input distribution, we assess their performance—previously trained on the Reddit TL;DR summarization task—on a distinct domain, namely news articles from the test split of the CNN/DailyMail dataset \citep{nallapati-etal-2016-abstractive}. For this evaluation, we employ the optimal sampling temperatures identified in the TL;DR setting (0 and 0.25). Table~\ref{tab:ood} shows the results. Consistent with earlier methodology, we measure GPT-4 win rates by comparing model outputs against ground-truth summaries, modifying our GPT-4 (C) prompt to substitute ``forum post'' with ``news article''. In this new context, DPO maintains a notable lead over PPO in performance. These findings suggest that DPO’s policies generalize at least as well as those produced by PPO, despite DPO not leveraging the unlabeled Reddit TL;DR prompts used in PPO training.

\subsection{Validating GPT-4 judgments with human judgments}
\label{sec:human-judgments}
We carry out a human evaluation to assess the trustworthiness of GPT-4’s judgments, drawing upon outcomes from the TL;DR summarization experiment and employing two types of GPT-4 prompts. The \textbf{GPT-4 (S)} (simple) prompt solicits a judgment about which summary better covers the key information in the post. By contrast, the \textbf{GPT-4 (C)} (concise) prompt includes an explicit criterion regarding conciseness, as our observations indicate that with the \textbf{GPT-4 (S)} prompt, GPT-4 demonstrates a preference for longer, more repetitive summaries than those favored by humans. The full text of the prompts can be found in Appendix~\ref{app:prompts}. We set up three head-to-head evaluations: the top-performing method (DPO, temperature 0.25), the lowest-performing method (PPO, temperature 1.0), and a method with intermediate performance (SFT, temperature 0.25), each tested against PPO with greedy sampling (its best temperature). These comparisons were selected to represent a wide range of summary qualities; for each, both prompt variants were used. Results indicate that for both prompt styles, GPT-4’s preferences align with human judgment to a degree comparable with inter-annotator human agreement—especially in DPO and PPO-1 cases, where multiple human annotations are available. As a general trend, the \textbf{GPT-4 (C)} prompt yields win rates more closely mirroring human decisions; accordingly, it is adopted for the principal results in Section~\ref{sec:dpo-real-datasets}. For a detailed description of the human study procedure, including the rater interface and the roster of human participants, refer to Appendix~\ref{app:human-study}.

\section{Discussion}
Preference-based learning offers a robust and scalable method for developing proficient and well-aligned language models. In this work, we have presented DPO, a streamlined framework that facilitates training language models using preferences, circumventing the need for reinforcement learning techniques. Instead of reformulating the preference learning challenge to fit within traditional RL frameworks and leveraging standard RL algorithms, DPO proposes a correspondence between language model policies and reward functions. This approach makes it possible to optimize language models to align with human preferences \textit{directly} using a basic cross-entropy objective, all without relying on reinforcement learning or sacrificing generality. Remarkably, with minimal hyperparameter adjustment, DPO achieves performance on par with or superior to prevailing RLHF methods, including those utilizing PPO, thereby significantly lowering the practical hurdles to training language models from human input.

\textbf{Limitations \& Future Work.} Our findings point to a number of critical directions for subsequent investigation. One question concerns the ability of the DPO policy to generalize to out-of-distribution data compared to models trained on explicit reward functions. Preliminary evidence indicates that DPO-trained models generalize at least as well as those optimized with PPO, yet a thorough empirical evaluation remains necessary. Further, it is worth examining whether strategies such as self-labeling using DPO-trained models can leverage unlabeled prompts with similar efficacy. Additional topics for inquiry include understanding how issues like reward over-optimization present themselves in the DPO framework, and whether the minor dip in performance observed in Figure~\ref{fig:dialogue-main}-right could be attributed to this phenomenon. Moreover, although our experiments include models containing up to 6B parameters, extending DPO to contemporary models of much greater scale is an especially promising line of future research. Regarding evaluation protocols, our study observes that GPT-4-derived win rates are prompt-dependent; exploring optimal approaches to obtain reliable assessments from automated evaluators would be valuable. Lastly, the DPO approach holds promise for applications beyond aligning language models with human preferences, including its use in training generative models across diverse modalities.

\end{document}